\chapter{FPGA System Methodology and Implementation}
\label{ch:fpga_implementation}

This chapter delineates the rigorous methodological approach and low-level RTL implementation used to realize the proposed AI-integrated ultra-HFT system. The central hypothesis that hardware-accelerated reinforcement learning (RL) can provide a superior risk-adjusted return compared to traditional heuristic-based models necessitates a hybrid architecture. We bifurcate the system into two distinct operational domains: the \textbf{Ultra-Low-Latency (ULL) Data Path} for deterministic execution, and the \textbf{Hybrid Control Plane} for asynchronous model management and telemetry.

\section{The Ultra-Low-Latency Data Path: "Pure-in-Gates" Philosophy}
The data path is architected using a "pure-in-gates" philosophy, ensuring that every operation on the critical path, from network ingress to order dispatch, is implemented as a combinatorial or pipelined synchronous logic circuit. This eliminates the non-determinism associated with soft-processor cores or operating system interrupts.

\subsection{Network Ingress: Cycle-Accurate Deterministic Parsing}
The ingress pipeline minimizes latency by processing network headers in parallel with the data stream arrival. The \texttt{rx\_parser} module implements a hierarchical **Finite State Machine (FSM)** that peels away encapsulation layers at line-rate. The FSM states are synchronized to the 322.26 MHz clock ($T_{cyc} \approx 3.1$ ns):
\begin{itemize}
    \item \textbf{IDLE/SOF:} The FSM waits for the Start of Frame (SOF) delimiter from the 10GbE MAC.
    \item \textbf{ETH\_EXTRACT (Cycle 0):} Consumes the 14-byte Ethernet header. It utilizes combinatorial comparators to validate the EtherType (0x0800) and the Destination MAC address.
    \item \textbf{IP\_EXTRACT (Cycle 1):} Extracts the 20-byte IPv4 header. A pipelined **checksum validator** computes the 16-bit one's complement sum in real-time, asserting an \texttt{error\_drop} signal if the integrity check fails.
    \item \textbf{UDP\_EXTRACT (Cycle 2):} Parses the 8-byte UDP header. The module performs a parallel lookup against a programmable **Port Whitelist** (stored in Registers/LUTs) to filter for specific multicast market data feeds (e.g., NASDAQ ITCH).
    \item \textbf{PAYLOAD\_STR:} Once validated, the FSM transitions to the payload streaming state, asserting the \texttt{TVALID} signal of the **AXI-Stream (AXI-S)** interface to the decoder.
\end{itemize}

\subsection{Feed Handling: Semantic Extraction and Bit-Slicing}
The \texttt{itch\_decoder} transforms the raw byte stream into structured market events. Due to the unaligned nature of ITCH 5.0 messages across 64-bit AXI-S words, the decoder employs a **512-bit Barrel Shifter**. This logic aligns the message header to the MSB of the internal bus, allowing for single-cycle field extraction:
\begin{equation}
    P_{raw} = \text{DataBus}[47:16], \quad Q_{raw} = \text{DataBus}[15:0]
\end{equation}
The decoder performs **Endianness swapping** (Big-to-Little) using a combinatorial wiring transpose, converting the wire-format integers into hardware-native formats for the Order Book.

\section{Core Innovation: The RL-Inference Systolic Array}
The pivotal contribution of this research is the \texttt{strat\_decide} module, which embeds the trained DRL policy (from Chapter \ref{ch:rl_strategy}) directly into the FPGA logic. To meet the 15ns inference target, we architected the neural network as a **fully unrolled systolic array**.

\subsection{DSP-Accelerated MAC Units}
The inference core utilizes **DSP48 slices** (specialized arithmetic hardware blocks in the FPGA) to perform high-speed Multiply-Accumulate (MAC) operations. For a layer with input vector $x$ and weight matrix $W$, the neuron output $y_i$ is computed as:
\begin{equation}
    y_i = \sigma \left( \sum_{j=1}^{N} w_{i,j} \cdot x_j + b_i \right)
\end{equation}
Variables in the hardware MAC:
\begin{itemize}
    \item $w_{i,j}, b_i$: **Quantized 16-bit fixed-point** weights and biases, loaded into the FPGA during the configuration phase via the Control Plane.
    \item $\sigma$: The \textbf{Activation Function}, implemented as a Piecewise Linear (PWL) approximation of the ReLU or Sigmoid function using LUTs (Look-Up Tables).
\end{itemize}

\subsection{Pipelined Inference Stages}
To maximize throughput, the RL Core is decomposed into a 4-stage pipeline:
\begin{enumerate}
    \item \textbf{Stage 1: Feature Normalization.} The raw BBO and OBI signals are scaled to the range $[-1, 1]$ using a bit-shift operator (dividing by powers of 2 for efficiency).
    \item \textbf{Stage 2: Hidden Layer 1.} Performs the first set of parallel MAC operations across 32 DSP slices.
    \item \textbf{Stage 3: Hidden Layer 2.} Processes the intermediate activations through a second, pruned weight matrix.
    \item \textbf{Stage 4: Softmax/Thresholding.} The final output layer determines the optimal action $a^*$. If $a^* > \text{threshold}$, a trade signal is emitted to the OUCH encoder.
\end{enumerate}

\section{Hardware Risk Management: Pre-Trade Gating}
To mitigate the risks of "hallucination" in probabilistic RL models, the outputs of the AI core are gated by the \texttt{risk\_gate} module. This deterministic safety envelope enforces hard constraints:
\begin{itemize}
    \item \textbf{Max Notional Exposure:} $Q_{total} \cdot P_{exec} < \text{Limit}_{notional}$. This check is performed using a 64-bit hardware multiplier.
    \item \textbf{Message Rate Limiter:} Implements a \textbf{Token Bucket} algorithm to prevent "Quote Stuffing." If the number of orders sent within a 1ms window exceeds the threshold, subsequent orders are suppressed.
    \item \textbf{Inventory Hard Limits:} Prevents the RL agent from exceeding $q_{max}$ or $q_{min}$, forcing a "Liquidation-Only" mode if limits are reached.
\end{itemize}

\section{Hardware-Software Co-Design via PCIe Gen4}
The system utilizes a **PCIe Gen4 x16** interface (throughput $\approx 31.5$ GB/s) to bridge the FPGA with the software management layer.
\begin{itemize}
    \item \textbf{AXI-Lite Control Interface:} The software uses memory-mapped I/O (MMIO) to write model weights and risk parameters into the FPGA's configuration registers. This allows for **Hot-Swapping** of the RL policy without stopping the hardware pipeline.
    \item \textbf{AXI-Memory Mapped (AXI-MM) Telemetry:} The FPGA writes high-resolution timestamps ($1$ ns precision) and execution logs into a shared DMA region in the host's **Huge-Page RAM**. This provides the "Godhood" level observability needed for performance validation.
\end{itemize}

\section{Timing Closure and Physical Constraints}
Achieving timing closure at 322.26 MHz on a modern Kintex UltraScale+ FPGA requires rigorous physical constraints. We utilize **Floorplanning** (P-blocks) to co-locate the MAC units and the Order Book logic, minimizing the wire delay ($T_{prop}$) between features and inference. The design is verified to have a **Positive Slack** of $> 0.2$ ns across all PVT (Process, Voltage, Temperature) corners, ensuring operational stability in high-density data centers.

\section{Conclusion}
The FPGA implementation detailed in this chapter represents a paradigm shift in HFT architecture. By synthesizing the speed of deterministic RTL with the intelligence of pipelined RL inference, we create a system capable of adapting to market toxicity at wire speed. The following chapters detail the software stack and the empirical results of this hardware-accelerated approach.
